\section{\personNovikov{}覆盖与\personNovikov{}环}
\label{sec:novikov-covering-ring}

设 $(L_0,L_1)$ 是 $(M,\omega)$ 中一对相互横截的紧拉格朗日子流形。
\footnote{译者注：原书在本句中先称 $L_0,L_1$ 连通，下一句却特意声明此时不假设
它们连通；下一页还先用基准道路选出各自的连通分支，才说从此可以假设二者连通。
译文依这段论证删去本句中提前出现的“连通”。}
需要指出，此刻我们并\emph{不}假设拉格朗日子流形是连通的。

考虑道路空间
\[
  \Omega=\Omega(L_0,L_1)
  =\{\ell:[0,1]\longrightarrow M
    \mid \ell(0)\in L_0,\ \ell(1)\in L_1\}.
\]
\footnote{译者注：原书把道路 $\ell$ 的值域写成未定义的 $P$；本节的环境辛流形
始终为 $M$，译文依此补正。}
\stfhIndex{I026}{Zuoyongyixingshi}{作用 1-形式（action one-form）}
在 $\Omega$ 上定义一个 $1$-形式（\emph{作用 $1$-形式}）$\alpha$：
\begin{equation}
  \alpha(\ell)(\xi)
  =\int_0^1\omega\bigl(\dot\ell(t),\xi(t)\bigr)\,dt,
\label{eq:lagrangian-action-one-form}
\end{equation}
其中 $\xi\in T_\ell\Omega$；由推论 3.8.4 可知，它是一个“\emph{闭} $1$-形式”。

$\Omega(L_0,L_1)$ 并不连通，但它有可数多个连通分支。以后研究\personFloer{}模空间时，
我们总是在 $\Omega(L_0,L_1)$ 的某个固定连通分支上工作。选取一条\emph{\termBasePath{}} $\ell_0$，便指定了这个分支；把相应分支记为
\[
  \Omega(L_0,L_1;\ell_0)\subset\Omega(L_0,L_1).
\]

基准道路 $\ell_0$ 的始点和终点
\[
  x_0=\ell_0(0)\in L_0,
  \qquad
  x_1=\ell_0(1)\in L_1
\]
自动分别从 $L_0,L_1$ 中选出一个连通分支。于是
$\Omega(L_0,L_1;\ell_0)$ 无非就是这两个相应连通分支之间的道路空间。
因此，除非另有说明，从现在起\emph{我们总是假设 $L_0,L_1$ 连通}。

\stfhIndex{I138}{Nuoweikefufugai}{诺维科夫覆盖（Novikov covering）}
接下来描述一个覆盖空间，称为 $\Omega(L_0,L_1)$ 的分支
$\Omega(L_0,L_1;\ell_0)$ 的\emph{\termNovikovCovering{}}。我们先描述
$\Omega(L_0,L_1;\ell_0)$ 的\termUniversalCoveringSpace{}。考虑所有满足下列
边界条件的对子 $(\ell,w)$：
\begin{align}
  w(0,\mathord\cdot)&=\ell_0,
  &w(1,\mathord\cdot)&=\ell,
  \label{eq:novikov-filling-path-ends}\\
  w(s,0)&\in L_0,
  &w(s,1)&\in L_1
  \quad\text{对所有 }s\in[0,1].
  \label{eq:novikov-filling-boundary}
\end{align}
把 $w$ 看成 $\Omega(L_0,L_1;\ell_0)$ 中一条从 $\ell_0$ 到 $\ell$ 的连续道路，
则 $\Omega(L_0,L_1;\ell_0)$ 的万有覆盖空间在 $\ell$ 上的纤维，可以由 $w$
相对于两端 $s=0,1$ 的道路同伦类集合表示。\footnote{译者注：原书此处写成
“from $\ell_0$ and $\ell$”；按万有覆盖的道路模型以及本句紧接的端点条件，应为
“从 $\ell_0$ 到 $\ell$”。}
把 $\Omega(L_0,L_1;\ell_0)$ 中满足
\eqref{eq:novikov-filling-path-ends}的 $w$ 的这个同伦类记为 $[\ell,w]$。

现在用另一个比同伦更弱的等价关系对上述万有覆盖空间取商，由此定义
$\Omega(L_0,L_1;\ell_0)$ 的一个较小覆盖空间。这类似于辛流形
$(M,\omega)$ 的可缩环路空间的\personNovikov{}覆盖空间；我们将在第 4 部分研究后者。

注意，若给定两个对子 $(\ell,w)$ 与 $(\ell,w')$，其公共道路
$\ell\in\Omega(L_0,L_1;\ell_0)$，则拼接
\[
  \overline w\#w':[0,1]\times[0,1]\longrightarrow M
\]
定义了一个环路
\[
  c:S^1\longrightarrow\Omega(L_0,L_1;\ell_0).
\]
可以把这个环路视为一个\termAnnularMap{}
\[
  C:S^1\times[0,1]\longrightarrow M,
\]
它满足边界条件
\begin{equation}
  C(s,0)\in L_0,
  \qquad
  C(s,1)\in L_1.
\label{eq:annular-map-lagrangian-boundary}
\end{equation}
显然，$C$ 的\termSymplecticArea{}记为
\begin{equation}
  I_\omega(C)=\int_C\omega,
\label{eq:annular-symplectic-area}
\end{equation}
它只依赖于满足\eqref{eq:annular-map-lagrangian-boundary}的 $C$ 的同伦类，
因而定义了一个同态，记为
\[
  I_\omega:
  \pi_1\bigl(\Omega(L_0,L_1;\ell_0)\bigr)
  \longrightarrow\mathbb R.
\]

接下来注意，对满足\eqref{eq:annular-map-lagrangian-boundary}的映射
$C:S^1\times[0,1]\to M$，可以关联一个辛丛对 $(E_C,\lambda_C)$：
\[
  E_C=C^*TM,
  \qquad
  \lambda_C=R_0^*TL_0\sqcup R_1^*TL_1,
\]
其中 $R_i:S^1\to L_i$ 是映射
\[
  R_i(s)=C(s,i),
  \qquad i=0,1.
\]
这使我们能够定义另一个同态
\[
  I_\mu:
  \pi_1\bigl(\Omega(L_0,L_1;\ell_0),\ell_0\bigr)
  \longrightarrow\mathbb Z;
  \qquad
  I_\mu(c)=\mu(E_C,\lambda_C),
\]
其中 $\mu(E_C,\lambda_C)$ 是丛对 $(E_C,\lambda_C)$ 的马斯洛夫指标。
\footnote{译者注：原书在 $I_\mu$ 的定义域中突然把已经固定的分支记号
$\Omega(L_0,L_1;\ell_0)$ 改写成未定义的 $\Omega_{\ell_0}(L_0,L_1)$；译文统一
沿用前者。}

利用 $I_\mu$ 与 $I_\omega$\footnote{译者注：原书在此把同态 $I_\omega$
称为“辛形式”；译文按上文已经给出的类型只写其记号。}，在所有满足
\eqref{eq:novikov-filling-path-ends}--\eqref{eq:novikov-filling-boundary}的对子
$(\ell,w)$ 之集合上定义一个等价关系 $\sim$。\footnote{译者注：原书误称这些对子
满足 (13.4.12)，但该式是环带映射 $C$ 的边界条件；对子 $(\ell,w)$ 的条件是
(13.4.10)--(13.4.11)，译文依此补正。}
对具有公共道路 $\ell$ 的两个这样的填充 $w,w'$，以
$\overline w\#w'$ 表示二者沿 $\ell$ 的拼接；它定义了
$\Omega(L_0,L_1;\ell_0)$ 中一条以 $\ell_0$ 为基点的环路。

\stfhIndex{I008}{00-gamma-equivalent}{$\Gamma$-等价（$\Gamma$-equivalent）}
\begin{definition}
\label{def:gamma-equivalent-fillings}
若
\begin{equation}
  I_\omega(\overline w\#w')=0
  =I_\mu(\overline w\#w'),
\label{eq:gamma-equivalence}
\end{equation}
则称 $(\ell,w)$ 与 $(\ell,w')$ 是 $\Gamma$-等价的，并记作
$(\ell,w)\sim(\ell,w')$。以
\stfhIndex{I006}{00-gamma-covering}{$\Gamma$-覆盖空间（$\Gamma$-covering space）}
\[
  \widetilde\Omega(L_0,L_1;\ell_0)
\]
表示等价类 $[\ell,w]$ 的集合，并称它为
$\Omega(L_0,L_1;\ell_0)$ 的\termGammaCoveringSpace{}。
\end{definition}

$\ell_0\in\Omega(L_0,L_1;\ell_0)$ 有一个到
$\widetilde\Omega(L_0,L_1;\ell_0)$ 的典范提升，即
\[
  [\ell_0,\widetilde\ell_0]
  \in\widetilde\Omega(L_0,L_1;\ell_0),
\]
其中 $\widetilde\ell_0$ 是映射
\[
  \widetilde\ell_0:[0,1]^2\longrightarrow M,
  \qquad
  \widetilde\ell_0(s,t)=\ell_0(t).
\]
这样，$\widetilde\Omega(L_0,L_1;\ell_0)$ 也有一个自然基点；我们在记号中省略它。

\stfhIndex{I002}{00-g-l0-l1-ell}{$G(L_0,L_1;\ell)$}
现在，以 $G(L_0,L_1;\ell)$ 表示作用在 $\pi^{-1}(\ell)$ 上的\termDeckTransformation{}群。不难看出，若两条不同道路 $\ell$ 相互同伦，
则相应的两个群之间有典范同构。于是两个同态 $I_\omega,I_\mu$
\footnote{译者注：原书在这里把此前明确定义为 homomorphisms 的
$I_\omega,I_\mu$ 称为 isomorphisms；译文依其类型统一为“同态”。}
下降为同态
\[
  \nu:G(L_0,L_1;\ell_0)\longrightarrow\mathbb R,
  \qquad
  d:G(L_0,L_1;\ell_0)\longrightarrow\mathbb Z,
\]
其定义为
\begin{equation}
  \nu(g)=\omega[g],
  \qquad
  d(g)=\mu(g).
\label{eq:deck-area-degree-homomorphisms}
\end{equation}

\begin{proposition}
\label{prop:novikov-deck-group-abelian}
群 $G(L_0,L_1;\ell_0)$ 是阿贝尔群。
\end{proposition}

\begin{proof}
由 $G(L_0,L_1;\ell_0)$ 的定义，映射
\[
  \nu\times d:G(L_0,L_1;\ell_0)
  \longrightarrow\mathbb R\times\mathbb Z
\]
是单射群同态。因为 $\mathbb R\times\mathbb Z$ 是阿贝尔群，所以
$G(L_0,L_1;\ell_0)$ 也是阿贝尔群。
\end{proof}

\stfhIndex{I151}{Zhouqiqun}{周期群（period group）}
\begin{definition}[周期群]
\label{def:lagrangian-period-group}
定义群 $G(L_0,L_1;\ell_0)$（或分支
$\widetilde\Omega(L_0,L_1;\ell_0)$）的\termPeriodGroup{}为像
\[
  \Gamma_\omega(L_0,L_1;\ell_0)
  =\{\omega[g]\mid g\in G(L_0,L_1;\ell_0)\}
  \subset\mathbb R.
\]
\end{definition}

现在定义与阿贝尔群
\[
  G(L_0,L_1;\ell_0)\subset\mathbb R\times\mathbb Z
\]
相伴的\termNovikovRing{}。先对一个交换环 $R$ 考虑\termGroupRing{}
\[
  R[G(L_0,L_1;\ell_0)].
\]
我们将用 $\nu$ 诱导的拓扑来定义环 $R[G(L_0,L_1;\ell_0)]$ 的一个完备化。

考虑\termFormalSum{}
\[
  \sum_{g\in G(L_0,L_1;\ell_0)}a_gg=: \sigma,
\]
并令
\[
  \operatorname{supp}(\sigma)
  =\{g\in G(L_0,L_1;\ell_0)\mid a_g\ne0\}.
\]

\stfhIndex{I143}{Nuoweikefuhuan}{诺维科夫环（Novikov ring）}
\begin{definition}
\label{def:lagrangian-novikov-ring}
定义与群 $G(L_0,L_1;\ell_0)$ 相伴的\termNovikovRing{}为
\footnote{译者注：原书在下式的有限性条件中使用自由变量 $C$，却没有写出量词；
由 $\nu$ 诱导的完备化条件要求该有限性对每个 $C\in\mathbb R$ 成立，译文将其明示。}
\[
  \Lambda(L_0,L_1;\ell_0)
  =\left\{\sigma\ \middle|\
    \#\bigl(\operatorname{supp}(\sigma)
      \cap\nu^{-1}((C,\infty))\bigr)<\infty
    \text{ 对每个 }C\in\mathbb R\right\},
\]
并在其上取显然的环结构。
\end{definition}

不难看出，$\Lambda(L_0,L_1;\ell_0)$ 构成一个\termGradedRModule{}：
\[
  \Lambda(L_0,L_1;\ell_0)
  =\bigoplus_{k\in\mathbb Z}\Lambda^k(L_0,L_1;\ell_0),
\]
其中 $\Lambda^k(L_0,L_1;\ell_0)$ 是由上述形式和 $\sigma$ 中满足
$\mu(g)=k$ 的 $g$ 所生成的 $R$-子模。
